% This LaTeX was auto-generated from MATLAB code.
% To make changes, update the MATLAB code and export to LaTeX again.

\documentclass{article}

\usepackage[utf8]{inputenc}
\usepackage[T1]{fontenc}
\usepackage{lmodern}
\usepackage{graphicx}
\usepackage{color}
\usepackage{hyperref}
\usepackage{amsmath}
\usepackage{amsfonts}
\usepackage{epstopdf}
\usepackage[table]{xcolor}
\usepackage{matlab}
\usepackage[paperheight=795pt,paperwidth=614pt,top=72pt,bottom=72pt,right=72pt,left=72pt,heightrounded]{geometry}

\sloppy
\epstopdfsetup{outdir=./}
\graphicspath{ {./lab4_media/} }

\matlabmultipletitles

\begin{document}

\matlabtitle{Задание 1}

\begin{par}
\begin{flushleft}
\textbackslash{}section\{Задача стабилизации с идеальным дифференцирующим звеном\}
\end{flushleft}
\end{par}

\begin{par}
\begin{flushleft}
Рассмотрим объект управления 2-го порядка, заданный дифференциальным уравнением
\end{flushleft}
\end{par}

\begin{par}
$$a_2 \ddot{y} +a_1 \dot{y} +a_0 y=u.~~(1)$$
\end{par}

\begin{par}
\begin{flushleft}
Придумаем такие коэффициенты $a_2$, $a_1$ и $a_0$, чтобы она содержала хотя бы один неустойчивый полюс. Зададимся ненулевым $\dot{y} (0)$ ($y(0)$ примем равным 0) и выполним моделирование свободного движения разомкнутой системы $y_{\textrm{раз}} (t)$.
\end{flushleft}
\end{par}

\begin{par}
\begin{flushleft}
Рассмотрим регулятор вида 
\end{flushleft}
\end{par}

\begin{par}
$$u=k_0 y+k_1 \dot{y} ~~(2)$$
\end{par}

\begin{par}
\begin{flushleft}
и построим структурную схему замкнутой системы, состоящей из объекта управления (1) и регулятора (2), в режиме стабилизации. Определим, при каких значениях параметров $k_0$ и $k_1$ замкнутая система будет устойчивой.
\end{flushleft}
\end{par}


\vspace{1em}
\begin{par}
\begin{flushleft}
Зададимся конкретными значениями параметров $k_0$ и $k_1$, обеспечивающими асимптотически устойчивую замкнутую систему, и выполним моделирования движения замкнутой системы $y_{\textrm{з}} (t)$ с начальными условиями, выбранными в рамках предыдущего моделирования. При моделировании в программной среде MATLAB/Simulink для получения производной $\dot{y} (t)$ используем блок Derivativ, основанный на использовании конечной разности с малым шагом $\Delta t$
\end{flushleft}
\end{par}

\begin{par}
$$\dot{y} (t)\approx \frac{y(t)-y(t-\Delta t)}{\Delta t}.$$
\end{par}

\begin{par}
\begin{flushleft}
Сравним результаты экспериментов и сделаем выводы.
\end{flushleft}
\end{par}


\vspace{1em}
\begin{par}
\begin{flushleft}
Найдем подходящие коэффициенты:
\end{flushleft}
\end{par}

\begin{par}
\begin{flushleft}
$a_2 \lambda^2 +a_1 \lambda +a_0 =0$, пусть $a_0 =-3,a_1 =-2,a_2 =1\Rightarrow (\lambda -3)(\lambda +1)=0,\lambda_1 =-1,\lambda_2 =3$.
\end{flushleft}
\end{par}

\begin{par}
\begin{flushleft}
Полюс $\lambda_2 =3$ неустойчив по корневому критерию. Возьмём $\dot{y} (0)=4$.
\end{flushleft}
\end{par}

\begin{par}
\begin{flushleft}
Построим структурную схему (см. рис. ???):
\end{flushleft}
\end{par}

\begin{par}
$$a_2 p^2 y+a_1 py+a_0 y=u\Rightarrow y=\frac{1}{a_2 p}\left(\frac{1}{p}(u-a_0 y)-a_1 y\right)$$
\end{par}

\begin{par}
\begin{flushleft}
???РИСУНОК???
\end{flushleft}
\end{par}

\begin{par}
\begin{flushleft}
Выполним моделирование (см. рис. ???).
\end{flushleft}
\end{par}

\begin{par}
\begin{flushleft}
???РИСУНОК???
\end{flushleft}
\end{par}

\begin{par}
\begin{flushleft}
Теперь построим структурную схему с регулятором $u=k_0 y+k_1 \dot{y}$ (см. рис. ???):
\end{flushleft}
\end{par}

\begin{par}
\begin{flushleft}
???РИСУНОК???
\end{flushleft}
\end{par}

\begin{par}
\begin{flushleft}
Определим при каких коэффициентах $k_i$ система будет устойчивой:
\end{flushleft}
\end{par}

\begin{par}
$$a_2 \ddot{y} +a_1 \dot{y} +a_0 y=k_0 y+k_1 \dot{y} \Rightarrow a_2 \ddot{y} +(a_1 -k_1 )\dot{y} +(a_0 -k_0 )y=0$$
\end{par}

\begin{par}
\begin{flushleft}
По теореме Стодолы (для случая систем не выше второго порядка, как у нас), система устойчива, если все коэффициенты характеристического полинома строго положительны, следовательно, для ассимптотической устойчивости:
\end{flushleft}
\end{par}

\begin{par}
$$a_2 >0,~~a_1 -k_1 >0,~~a_0 -k_0 >0\Rightarrow k_1 <-2,~~k_0 <-3$$
\end{par}

\begin{par}
\begin{flushleft}
Возьмём $k_0 =-5,k_1 =-4$. Выполним моделирование (см. рис. ???).
\end{flushleft}
\end{par}

\begin{par}
\begin{flushleft}
???РИСУНОК???
\end{flushleft}
\end{par}

\begin{par}
\begin{flushleft}
\textbf{Выводы}
\end{flushleft}
\end{par}

\begin{par}
\begin{flushleft}
 В данном задании мы сначала построили неустойчивую систему (при свободном движении), после чего воспользовались теоремой Стодолы для нахождения коэффициентов для регулятора, которые могут сделать замкнутую систему устойчивой и выполнили моделирования для системы с регулятором. Результаты сошлись с теоретическими ожиданиями. Это довольно наглядная и простая демонстрация того, как работает простейший регулятор.
\end{flushleft}
\end{par}

\begin{matlabcode}
a_0 = -3;
a_1 = -2;
a_2 = 1;

simtime = 3;
plot_1 = sim("lab4_1.slx", simtime);
plot(plot_1.y.Time, plot_1.y.Data, LineWidth=2, LineStyle="-", ...
            DisplayName=sprintf('$y(0)=%d, \\dot{y}(0)=%d$', 0, 4))
grid on;
legend(Location="best", FontSize=14, Interpreter="latex");
xlabel("$t, c$", Interpreter="latex");
ylabel("$y(t)$", Interpreter="latex");
set(gcf, Position=[100, 100, 1200, 600]);   
set(gca, FontSize=14);
exportgraphics(gcf, "images/plot_1.png", Resolution=200)
k_0 = -5;
k_1 = -4;

simtime = 10;
plot_1 = sim("lab4_1z.slx", simtime);
plot(plot_1.y.Time, plot_1.y.Data, LineWidth=2, color="red", LineStyle="-", ...
            DisplayName=sprintf('$y(0)=%d, \\dot{y}(0)=%d$', 0, 4))
grid on;
legend(Location="best", FontSize=14, Interpreter="latex");
xlabel("$t, c$", Interpreter="latex");
ylabel("$y(t)$", Interpreter="latex");
set(gcf, Position=[100, 100, 1200, 600]);   
set(gca, FontSize=14);
exportgraphics(gcf, "images/plot_1z.png", Resolution=200)
\end{matlabcode}


\matlabtitle{Задание 2}

\begin{par}
\begin{flushleft}
\textbackslash{}section\{Задача стабилизации с реальным дифференцирующим звеном\}
\end{flushleft}
\end{par}

\begin{par}
\begin{flushleft}
Модифицируем замкнутую систему из \textbf{Задания 1}, заменив аппроксимацию производной $\dot{y} (t)$ на передаточную функцию вида
\end{flushleft}
\end{par}

\begin{par}
$$W_{\textrm{р.дифф.}} (p)=\frac{p}{Tp+1}.~~(3)$$
\end{par}

\begin{par}
\begin{flushleft}
Определим аналитически критические значения параметра $T$, при которых система становится неустойчивой для выбранных ранее $k_0$ и $k_1$.
\end{flushleft}
\end{par}

\begin{par}
\begin{flushleft}
Осуществим аналогичное \textbf{Заданию 1} моделирование движения замкнутой системы $y_{\textrm{з}} (t)$ для нескольких различных значений параметра $T$, соответствующих устойчивой
\end{flushleft}
\end{par}

\begin{par}
\begin{flushleft}
системе. Приведем соответствующие графики выхода $y(t)$ и сопоставим их между собой и с результатом моделирования замкнутой системы из \textbf{Задания 1}. Сделаем выводы.
\end{flushleft}
\end{par}


\vspace{1em}
\begin{par}
\begin{flushleft}
Модифицируем схему (см. рис. ???).
\end{flushleft}
\end{par}

\begin{par}
\begin{flushleft}
???РИСУНОК???
\end{flushleft}
\end{par}

\begin{par}
\begin{flushleft}
Найдём значения $T$, при которых система становится неустойчивой:
\end{flushleft}
\end{par}

\begin{par}
$$\begin{array}{l}
a_2 p^2 y+a_1 py+a_0 y=k_0 y+k_1 \frac{p}{Tp+1}y,\\
a_2 Tp^3 y+a_2 p^2 y+a_1 Tp^2 y+(a_1 -k_1 -k_0 T)py+a_0 Tpy+(a_0 -k_0 )y=0,\\
Tp^3 y+(1-2T)p^2 y-(2+k_1 +k_0 T+3T)py-(3+k_0 )y=0,\\
T\dddot{y} +(1-2T)\ddot{y} +2(1+T)\dot{y} +2y=0
\end{array}$$
\end{par}

\begin{par}
\begin{flushleft}
Система однозначно неустойчива при отрицательных коэффициентах, то есть при: $T<0,~~T>\frac{1}{2}$.
\end{flushleft}
\end{par}

\begin{par}
\begin{flushleft}
Воспользуемся критерием Гурвица для определения областей устойчивости. Матрица Гурвица:
\end{flushleft}
\end{par}

\begin{par}
$$A=\left(\begin{array}{ccc}
1-2T & 2 & 0\\
T & 2(1+T) & 0\\
0 & 1-2T & 2
\end{array}\right)$$
\end{par}

\begin{par}
\begin{flushleft}
Ведущие угловые миноры матрицы Гурвица:
\end{flushleft}
\end{par}

\begin{par}
$$\begin{array}{l}
\Delta_1 =1-2T,\\
\Delta_2 =2-4T-4T^2 =-4\left(T+\frac{1+\sqrt{3}}{2}\right)\left(T+\frac{1-\sqrt{3}}{2}\right),\\
\Delta_3 =2\Delta_2 =4-8T-8T^2 =-8\left(T+\frac{1+\sqrt{3}}{2}\right)\left(T+\frac{1-\sqrt{3}}{2}\right)
\end{array}$$
\end{par}

\begin{par}
\begin{flushleft}
Следовательно, система ассимптотически устойчива при: $T\in \left(0,\frac{1+\sqrt{3}}{2}\right)$.
\end{flushleft}
\end{par}

\begin{par}
\begin{flushleft}
Выполним моделирование при различных $T$, соответствущих устойчивости система (см. рис. ???).
\end{flushleft}
\end{par}

\begin{par}
\begin{flushleft}
???РИСУНОК???
\end{flushleft}
\end{par}

\begin{par}
\begin{flushleft}
\textbf{Выводы}
\end{flushleft}
\end{par}

\begin{par}
\begin{flushleft}
В данном задании мы заменили прямое численное дифференцирование на передаточную функцию с параметром $T$ и выполнили моделирование для различных его значений, соответствующих устойчивой системе. Это временной параметр, который отражает скорость реакции системы не воздействие. При меньших значениях $T$ система быстрее приходит к установившемуся значению, однако, чем ближе $T$ к нулю, тем ближе мы и к физической нереализуемости системы, так как при $T=0$ числитель$W_{\textrm{р.дифф.}} (p)=\frac{p}{Tp+1}$ становится больше знаменателя. При увеличении $T$ увеличивается колебательность при переходном процессе и он длится дольше. При $T$, близком к нулю, мы видим процесс, очень похожий на использование дифференцирования.
\end{flushleft}
\end{par}

\begin{matlabcode}
T_arr = [0.3, 0.2, 0.1, 0.01];
simtime = 20;
for i = 1:length(T_arr)
    T = T_arr(i);
    plot_2 = sim("lab4_2z.slx", simtime);
    plot(plot_2.y.Time, plot_2.y.Data, LineWidth=3, LineStyle="-", ...
                DisplayName=sprintf('Transfer Fcn, $T=%.2f$', T))
    hold on;
end
plot_1 = sim("lab4_1z.slx", simtime);
plot(plot_1.y.Time, plot_1.y.Data, LineWidth=4, color="red", LineStyle="--", ...
            DisplayName='Derivative')
hold on;

grid on;
legend(Location="best", FontSize=14, Interpreter="latex");
xlabel("$t, c$", Interpreter="latex");
ylabel("$y(t)$", Interpreter="latex");
set(gcf, Position=[100, 100, 1200, 600]);   
set(gca, FontSize=14);
hold off;
exportgraphics(gcf, "images/plot_2z.png", Resolution=200)
\end{matlabcode}


\matlabtitle{Задание 3}

\begin{par}
\begin{flushleft}
\textbackslash{}section\{Задача слежения для системы с астатизмом нулевого порядка (П-регулятор)\}
\end{flushleft}
\end{par}

\begin{par}
\begin{flushleft}
Рассмотрим замкнутую систему, заданную структурной схемой на рисунке ???, где
\end{flushleft}
\end{par}

\begin{par}
$$H(s)=k$$
\end{par}

\begin{par}
\begin{flushleft}
— пропорциональный регулятор. Варианты передаточной функции объекта управления $W(s)$ и характеристики задающего воздействия $g(t)$ приведены в \textbf{Таблице 1 }инструкции по выполнению лабораторной работы. Зададимся не менее чем тремя значениями параметра $k$, которые бы обеспечивали асимптотически устойчивую замкнутую систему.
\end{flushleft}
\end{par}

\begin{par}
\begin{flushleft}
Исследуем стационарный режим работы при $g(t)=A$: выполним моделирование для выбранных значений $k$, а также аналитически определим предельное значение ошибки $e_{уст}$ для каждого значения $k$. Сопоставим результаты между собой и сделаем выводы о влиянии величины $k$.
\end{flushleft}
\end{par}

\begin{par}
\begin{flushleft}
Исследуем режим движения с постоянной скоростью при $g(t)=Vt$: выполним моделирование для выбранных значений $k$. Сопоставим результаты между собой и сделаем выводы о влиянии величины $k$.
\end{flushleft}
\end{par}

\begin{par}
\begin{flushleft}
???РИСУНОК???
\end{flushleft}
\end{par}

\begin{par}
\begin{flushleft}
Возьмём из таблицы для варианта 27 передаточную функцию и параметры управляющего воздействия:
\end{flushleft}
\end{par}

\begin{par}
$$W(s)=\frac{1}{4s^2 +0.6s+1},~~A=3,~~Vt=2.5t$$
\end{par}

\begin{par}
\begin{flushleft}
Передаточная функция замкнутой системы: $W_{g\to y} (s)=\frac{W(s)W_{\textrm{рег}} (s)}{1+W(s)W_{\textrm{рег}} (s)}$, где $W(s)W_{\textrm{рег}} (s)$ - передаточная функция разомкнутой системы.
\end{flushleft}
\end{par}

\begin{par}
$$W_{g\to y} (s)=\frac{W(s)W_{\textrm{рег}} (s)}{1+W(s)W_{\textrm{рег}} (s)}=\frac{\frac{k}{4s^2 +0.6s+1}}{1+\frac{k}{4s^2 +0.6s+1}}=\frac{k}{4s^2 +0.6s+1+k}$$
\end{par}

\begin{par}
\begin{flushleft}
Найдем такие $k$, при которых замкнутая система является ассимптотически устойчивой.
\end{flushleft}
\end{par}

\begin{par}
$$\begin{array}{l}
4s^2 +0.6s+1+k=0,~~s_{1,2} =\frac{-0.6\pm \sqrt{0.36-4\cdot 4(1+k)}}{8}=\frac{-0.6\pm \sqrt{-15.64-16k}}{8}=\\
=-0.075\pm \frac{\sqrt{-(15.64+16k)}}{8}\Rightarrow \textrm{Re}\left\lbrack \sqrt{-(15.64+16k)}\right\rbrack <0.6\Rightarrow k>-1
\end{array}$$
\end{par}

\begin{par}
\begin{flushleft}
При $k>-1$ полюса передаточной функции будут отрицательными и мы получаем нужную устойчивость.
\end{flushleft}
\end{par}

\begin{par}
\begin{flushleft}
Зададимся значениями $k=\lbrace 0.5,2,29\rbrace$.
\end{flushleft}
\end{par}

\begin{par}
\begin{flushleft}
Проведем исследование для двух режимов работы.
\end{flushleft}
\end{par}

\begin{par}
\begin{flushleft}
\textbf{Стационарный: }$G(s)=\frac{A}{s}$.
\end{flushleft}
\end{par}

\begin{par}
\begin{flushleft}
Аналитически определим предельное значение ошибки $e_{уст}$ для каждого значения $k$:
\end{flushleft}
\end{par}

\begin{par}
\begin{flushleft}
Ошибка слежения:
\end{flushleft}
\end{par}

\begin{par}
$$E(s)=\frac{1}{1+W(s)W_{\textrm{рег}} (s)}G(s)=\frac{1}{1+\frac{k}{4s^2 +0.6s+1}}\frac{A}{s}=\frac{A(4s^2 +0.6s+1)}{s(4s^2 +0.6s+1+k)}$$
\end{par}

\begin{par}
\begin{flushleft}
Согласно теореме о конечном значении, если все полюса $s\cdot E(s)$ имеют отрицательную вещественную часть, $\lim_{t\to +\infty } y(t)=\lim_{s\to 0} sY(s)$.
\end{flushleft}
\end{par}

\begin{par}
\begin{flushleft}
$s\cdot E(s)=\frac{A(4s^2 +0.6s+1)}{4s^2 +0.6s+1+k}$, полюса данной передаточной функции отрицательны, так как ранее мы уже подбирали диапазон $k$, соответствующий отрицательным полюсам передаточной функции с таким же знаменателем. Можем применнить теорему о конечном значении и найти установившуюся ошибку:
\end{flushleft}
\end{par}

\begin{par}
$$e_{\textrm{уст}} =\lim_{s\to 0} s\cdot E(s)=\lim_{s\to 0} \frac{A(4s^2 +0.6s+1)}{4s^2 +0.6s+1+k}=\frac{A(4\cdot 0^2 +0.6\cdot 0+1)}{4\cdot 0^2 +0.6\cdot 0+1+k}=\frac{A}{1+k}$$
\end{par}

\begin{par}
\begin{flushleft}
Рассчитаем для каждого $k$:
\end{flushleft}
\end{par}

\begin{itemize}
\setlength{\itemsep}{-1ex}
   \item{\begin{flushleft} $\displaystyle k=0.5,~~e_{\textrm{уст}} =\frac{3}{1+0.5}=2$ \end{flushleft}}
   \item{\begin{flushleft} $\displaystyle k=2,~~e_{\textrm{уст}} =\frac{3}{1+2}=1$ \end{flushleft}}
   \item{\begin{flushleft} $\displaystyle k=29,~~e_{\textrm{уст}} =\frac{3}{1+29}=0.1$ \end{flushleft}}
\end{itemize}

\begin{par}
\begin{flushleft}
Выполним моделирование. Структурная схема представлена на рис. ???.
\end{flushleft}
\end{par}

\begin{par}
\begin{flushleft}
???РИСУНОК???
\end{flushleft}
\end{par}

\begin{par}
\begin{flushleft}
Графики входа и выходов на рис. ???.
\end{flushleft}
\end{par}

\begin{par}
\begin{flushleft}
???РИСУНОК???
\end{flushleft}
\end{par}

\begin{par}
\begin{flushleft}
Графики ошибки на рис. ???.
\end{flushleft}
\end{par}

\begin{par}
\begin{flushleft}
???РИСУНОК???
\end{flushleft}
\end{par}

\begin{par}
\begin{flushleft}
Как мы видим из полученных графиков, при увеличении $k$ уменьшается установившаяся ошибка, но растёт перерегулирование и колебательность системы.
\end{flushleft}
\end{par}

\begin{par}
\begin{flushleft}
\textbf{Движение с постоянной скоростью: }$G(s)=\frac{V}{s^2 }$.
\end{flushleft}
\end{par}

\begin{par}
\begin{flushleft}
Выполним моделирование. Структурная схема представлена на рис. ???.
\end{flushleft}
\end{par}

\begin{par}
\begin{flushleft}
???РИСУНОК???
\end{flushleft}
\end{par}

\begin{par}
\begin{flushleft}
Графики входа и выходов на рис. ???.
\end{flushleft}
\end{par}

\begin{par}
\begin{flushleft}
???РИСУНОК???
\end{flushleft}
\end{par}

\begin{par}
\begin{flushleft}
Графики ошибки на рис. ???.
\end{flushleft}
\end{par}

\begin{par}
\begin{flushleft}
???РИСУНОК???
\end{flushleft}
\end{par}

\begin{par}
\begin{flushleft}
Из графиков видно, что следить за таким входным воздействием у пропорционального регулятора уже не получается - ошибка растет линейно (с точностью аппроксимации функции ошибки). При увеличении $k$ ошибка накапливается медленнее, но при $t\to \infty$ она всё равно уйдет в бесконечность из-за недостающего порядка астатизма.
\end{flushleft}
\end{par}

\begin{par}
\begin{flushleft}
\textbf{Выводы}
\end{flushleft}
\end{par}

\begin{par}
\begin{flushleft}
В данном задании мы исследовали пропорциональный регулятор. Для анализа системы было удобно использовать теорему о конечном значении.
\end{flushleft}
\end{par}

\begin{par}
\begin{flushleft}
При исследовании стационарного режима были аналитически получены значения установившейся ошибки, которые подтверждаются результатами моделирования. Увеличение коэффициента усиления $k$ приводит к уменьшению установившейся ошибки слежения. Однако вместе с этим заметно возрастает перерегулирование и колебательность переходного процесса, что ухудшает качество регулирования и может ограничивать максимально допустимое значение $k$ в практических условиях.
\end{flushleft}
\end{par}

\begin{par}
\begin{flushleft}
При исследовании режима движения с постоянной скоростью показано, что пропорциональный регулятор не обеспечивает нулевой или фиксированной установившейся ошибки при линейно растущем задающем воздействии. Ошибка увеличивается со временем, а увеличение коэффициента усиления лишь уменьшает скорость её нарастания, но не устраняет её полностью. Это объясняется тем, что система с П-регулятором обладает \textbf{нулевым порядком астатизма}, следовательно, она способна компенсировать только ступенчатые воздействия (с точностью до фиксированной установившейся ошибки), но не линейно (и быстрее) растущие.
\end{flushleft}
\end{par}

\begin{par}
\begin{flushleft}
Таким образом, использование пропорционального регулятора является недостаточным для задач слежения за возрастающим сигналом, но можно использовать для слежения за константным там, где можно допустить небольшую установившуюся ошибку.
\end{flushleft}
\end{par}

\begin{matlabcode}
k_arr = [ 29, 2, 0.5 ];
A = 3;
V = 2.5;

simtime = 80;
figure;
set(gcf,'AutoResizeChildren','off');
hold on;

for i = 1:length(k_arr)
    k = k_arr(i);
    plot_3 = sim("lab4_31.slx", simtime);
    plot(plot_3.y.Time, plot_3.y.Data, LineWidth=3, LineStyle="-", ...
                DisplayName=sprintf('$y(t), k=%.1f$', k))
    hold on;
end
plot_31 = sim("lab4_31.slx", simtime);
plot(plot_31.g.Time, plot_31.g.Data, LineWidth=3, color="red", LineStyle="--", ...
            DisplayName='$g(t)=3$')
hold on;

grid on;
legend(Location="best", FontSize=14, Interpreter="latex");
xlabel("$t, c$", Interpreter="latex");
ylabel("$f(t)$", Interpreter="latex");
set(gcf, Position=[100, 100, 1200, 600]);   
set(gca, FontSize=14);

drawnow;
exportgraphics(gcf, "images/plot_31.png", Resolution=200)
hold off;
e_inf = [0.1, 1, 2]; 

figure;
set(gcf,'AutoResizeChildren','off');
hold on;
colors = hsv(length(k_arr));

for i = 1:length(k_arr)
    k = k_arr(i);
    plot_3 = sim("lab4_31.slx", simtime);
    plot(plot_3.e.Time, plot_3.e.Data, LineWidth=3, LineStyle="-", ...
                DisplayName=sprintf('$e(t), k=%.1f$', k), Color=colors(i,:))
    darkColor = colors(i,:) * (1 - 0.5);

    yline(e_inf(i), '--', Color=darkColor, LineWidth=2, ...
        DisplayName=sprintf('$e_{\\infty}, k=%.1f$', k));

    hold on;
end

grid on;
legend(Location="best", FontSize=14, Interpreter="latex");
xlabel("$t, c$", Interpreter="latex");
ylabel("$e(t)$", Interpreter="latex");
set(gcf, Position=[150, 150, 1200, 600]);
set(gca, FontSize=14);

drawnow;
exportgraphics(gcf, "images/plot_31_err.png", Resolution=200)
hold off;
\end{matlabcode}


\begin{matlabcode}
simtime = 50;
figure;
set(gcf,'AutoResizeChildren','off');
hold on;

for i = 1:length(k_arr)
    k = k_arr(i);
    plot_3 = sim("lab4_32.slx", simtime);
    plot(plot_3.y.Time, plot_3.y.Data, LineWidth=2, LineStyle="-", ...
                DisplayName=sprintf('$y(t), k=%.1f$', k))
    hold on;
end
plot_31 = sim("lab4_32.slx", simtime);
plot(plot_31.g.Time, plot_31.g.Data, LineWidth=2, color="red", LineStyle="--", ...
            DisplayName='$g(t)=2.5t$')
hold on;

grid on;
legend(Location="best", FontSize=14, Interpreter="latex");
xlabel("$t, c$", Interpreter="latex");
ylabel("$f(t)$", Interpreter="latex");
set(gcf, Position=[100, 100, 1200, 600]);   
set(gca, FontSize=14);

drawnow;
exportgraphics(gcf, "images/plot_32.png", Resolution=200)
hold off;
figure;
set(gcf,'AutoResizeChildren','off');
hold on;

for i = 1:length(k_arr)
    k = k_arr(i);
    plot_3 = sim("lab4_32.slx", simtime);
    plot(plot_3.e.Time, plot_3.e.Data, LineWidth=2, LineStyle="-", ...
                DisplayName=sprintf('$e(t), k=%.1f$', k))

    hold on;
end

grid on;
legend(Location="best", FontSize=14, Interpreter="latex");
xlabel("$t, c$", Interpreter="latex");
ylabel("$e(t)$", Interpreter="latex");
set(gcf, Position=[150, 150, 1200, 600]);
set(gca, FontSize=14);

drawnow;
exportgraphics(gcf, "images/plot_32_err.png", Resolution=200)
hold off;
\end{matlabcode}


\matlabtitle{Задание 4}

\begin{par}
\begin{flushleft}
\textbackslash{}section\{Задача слежения для системы с астатизмом первого порядка (И-регулятор)\}
\end{flushleft}
\end{par}

\begin{par}
\begin{flushleft}
Рассмотрим замкнутую систему, заданную структурной схемой на рисунке ???ВЫШЕ???, где
\end{flushleft}
\end{par}

\begin{par}
$$H(s)=\frac{k}{s}$$
\end{par}

\begin{par}
\begin{flushleft}
— интегральный регулятор. Варианты передаточной функции объекта управления $W(s)$ и характеристики задающего воздействия $g(t)$ приведены в  \textbf{Таблице 1 }инструкции по выполнению лабораторной работы. Зададимся не менее чем тремя значениями параметра $k$, которые бы обеспечивали асимптотически устойчивую замкнутую систему.
\end{flushleft}
\end{par}

\begin{par}
\begin{flushleft}
Исследуем стационарный режим работы при $g(t)=A$: выполним моделирование для выбранных значений $k$, а также аналитически определим предельное значение ошибки $e$ для каждого значения $k$. Сопоставим результаты между собой и сделаем выводы.
\end{flushleft}
\end{par}

\begin{par}
\begin{flushleft}
Исследуем режим движения с постоянной скоростью при $g(t)=Vt$: выполним моделирование для выбранных значений $k$, а также аналитически определим предельное значение ошибки $e$ для каждого значения $k$. Сопоставим результаты между собой и сделаем выводы.
\end{flushleft}
\end{par}

\begin{par}
\begin{flushleft}
Исследуем режим движения с постоянным ускорением при $g(t)=\frac{at^2 }{2}$: выполним моделирование для выбранных значений $k$. Сопоставим результаты между собой и сделаем выводы.
\end{flushleft}
\end{par}


\vspace{1em}
\begin{par}
\begin{flushleft}
Возьмём из таблицы для варианта 27 передаточную функцию и параметры управляющего воздействия:
\end{flushleft}
\end{par}

\begin{par}
$$W(s)=\frac{1}{4s^2 +0.6s+1},~~A=3,~~V=2.5,~~\frac{a}{2}=0.4$$
\end{par}

\begin{par}
\begin{flushleft}
Передаточная функция замкнутой системы:
\end{flushleft}
\end{par}

\begin{par}
$$W_{g\to y} (s)=\frac{W(s)H(s)}{1+W(s)H(s)}=\frac{\frac{k}{s}\cdot \frac{1}{4s^2 +0.6s+1}}{1+\frac{k}{s}\cdot \frac{1}{4s^2 +0.6s+1}}=\frac{k}{4s^3 +0.6s^2 +s+k}$$
\end{par}

\begin{par}
\begin{flushleft}
Найдем такие $k$, при которых замкнутая система является ассимптотически устойчивой.
\end{flushleft}
\end{par}

\begin{par}
\begin{flushleft}
$4s^3 +0.6s^2 +s+k=0$, воспользуемся критерием Гурвица, матрица Гурвица:
\end{flushleft}
\end{par}

\begin{par}
$$A=\left(\begin{array}{ccc}
0.6 & k & 0\\
4 & 1 & 0\\
0 & 0.6 & k
\end{array}\right)$$
\end{par}

\begin{par}
\begin{flushleft}
Ведущие угловые миноры матрицы Гурвица: $\Delta_1 =0.6,~~\Delta_2 =0.6-4k,~~\Delta_3 =k\cdot \Delta_2 =k(0.6-4k)$
\end{flushleft}
\end{par}

\begin{par}
\begin{flushleft}
Следовательно, система ассимптотически устойчива при: $k\in \left(0,0.15)\right.$.
\end{flushleft}
\end{par}

\begin{par}
\begin{flushleft}
Зададимся значениями $k=\lbrace 0.125,0.07,0.02\rbrace$.
\end{flushleft}
\end{par}

\begin{par}
\begin{flushleft}
Проведем исследование для трёх режимов работы.
\end{flushleft}
\end{par}

\begin{par}
\begin{flushleft}
\textbf{Стационарный: }$G(s)=\frac{A}{s}$.
\end{flushleft}
\end{par}

\begin{par}
\begin{flushleft}
Аналитически определим предельное значение ошибки $e_{уст}$ для каждого значения $k$:
\end{flushleft}
\end{par}

\begin{par}
\begin{flushleft}
Ошибка слежения:
\end{flushleft}
\end{par}

\begin{par}
$$E(s)=\frac{1}{1+W(s)H(s)}G(s)=\frac{s(4s^2 +0.6s+1)}{4s^3 +0.6s^2 +s+k}\cdot \frac{A}{s}=\frac{12s^2 +1.8s+3}{4s^3 +0.6s^2 +s+k}$$
\end{par}

\begin{par}
\begin{flushleft}
Согласно теореме о конечном значении, если все полюса имеют отрицательную вещественную часть, $\lim_{t\to +\infty } y(t)=\lim_{s\to 0} sY(s)$.
\end{flushleft}
\end{par}

\begin{par}
\begin{flushleft}
$s\cdot E(s)=\frac{s(12s^2 +1.8s+3)}{4s^3 +0.6s^2 +s+k}$, полюса данной передаточной функции отрицательны, так как ранее мы уже подбирали диапазон $k$, соответствующий отрицательным полюсам передаточной функции с таким же знаменателем. Можем применнить теорему о конечном значении и найти установившуюся ошибку:
\end{flushleft}
\end{par}

\begin{par}
$$e_{\textrm{уст}} =\lim_{s\to 0} s\cdot E(s)=\lim_{s\to 0} \frac{s(12s^2 +1.8s+3)}{4s^3 +0.6s^2 +s+k}=\frac{0\cdot (12\cdot 0^2 +1.8\cdot 0+3)}{4\cdot 0^3 +0.6\cdot 0^2 +s+k}=0$$
\end{par}

\begin{par}
\begin{flushleft}
Для любых $k$ установившаяся ошибка нулевая.
\end{flushleft}
\end{par}

\begin{par}
\begin{flushleft}
Выполним моделирование. Структурная схема представлена на рис. ???.
\end{flushleft}
\end{par}

\begin{par}
\begin{flushleft}
???РИСУНОК???
\end{flushleft}
\end{par}

\begin{par}
\begin{flushleft}
Графики входа и выходов на рис. ???.
\end{flushleft}
\end{par}

\begin{par}
\begin{flushleft}
???РИСУНОК???
\end{flushleft}
\end{par}

\begin{par}
\begin{flushleft}
Графики ошибки на рис. ???.
\end{flushleft}
\end{par}

\begin{par}
\begin{flushleft}
???РИСУНОК???
\end{flushleft}
\end{par}

\begin{par}
\begin{flushleft}
Видим, что $y(t)$ сходится к задающему воздействию при любых $k$, но при повышении $k$ возрастает колебательность переходного процесса. При уменьшении $k$ увеличивается время переходного процесса. Золотая середина из выбранных $k$ это $0.07$ - при нём колебательность незначительная и время переходного процесса наименьшее. У слишком маленьких значений $k$ время переходного процесса увеличивается из-за высокой колебательности.
\end{flushleft}
\end{par}

\begin{par}
\begin{flushleft}
\textbf{Движение с постоянной скоростью: }$G(s)=\frac{V}{s^2 }$.
\end{flushleft}
\end{par}

\begin{par}
\begin{flushleft}
Аналитически определим предельное значение ошибки $e_{уст}$ для каждого значения $k$:
\end{flushleft}
\end{par}

\begin{par}
\begin{flushleft}
Ошибка слежения:
\end{flushleft}
\end{par}

\begin{par}
$$E(s)=\frac{1}{1+W(s)H(s)}G(s)=\frac{s(4s^2 +0.6s+1)}{4s^3 +0.6s^2 +s+k}\cdot \frac{V}{s^2 }=\frac{V(4s^2 +0.6s+1)}{s(4s^3 +0.6s^2 +s+k)}$$
\end{par}

\begin{par}
\begin{flushleft}
Согласно теореме о конечном значении, если все полюса $s\cdot E(s)$ имеют отрицательную вещественную часть, $\lim_{t\to +\infty } y(t)=\lim_{s\to 0} sY(s)$.
\end{flushleft}
\end{par}

\begin{par}
\begin{flushleft}
$s\cdot E(s)=\frac{V(4s^2 +0.6s+1)}{4s^3 +0.6s^2 +s+k}$, полюса данной передаточной функции отрицательны, так как ранее мы уже подбирали диапазон $k$, соответствующий отрицательным полюсам передаточной функции с таким же знаменателем. Можем применнить теорему о конечном значении и найти установившуюся ошибку:
\end{flushleft}
\end{par}

\begin{par}
$$e_{\textrm{уст}} =\lim_{s\to 0} s\cdot E(s)=\lim_{s\to 0} \frac{V(4s^2 +0.6s+1)}{4s^3 +0.6s^2 +s+k}=\frac{V(4\cdot 0^2 +0.6\cdot 0+1)}{4\cdot 0^3 +0.6\cdot 0^2 +0+k}=\frac{V}{k}$$
\end{par}

\begin{par}
\begin{flushleft}
Рассчитаем для каждого $k$:
\end{flushleft}
\end{par}

\begin{itemize}
\setlength{\itemsep}{-1ex}
   \item{\begin{flushleft} $\displaystyle k=0.02,~~e_{\textrm{уст}} =\frac{2.5}{0.02}=125$ \end{flushleft}}
   \item{\begin{flushleft} $\displaystyle k=0.07,~~e_{\textrm{уст}} =\frac{2.5}{0.07}\approx 35.7$ \end{flushleft}}
   \item{\begin{flushleft} $\displaystyle k=0.125,~~e_{\textrm{уст}} =\frac{2.5}{0.125}=20$ \end{flushleft}}
\end{itemize}

\begin{par}
\begin{flushleft}
Выполним моделирование. Структурная схема представлена на рис. ???.
\end{flushleft}
\end{par}

\begin{par}
\begin{flushleft}
???РИСУНОК???
\end{flushleft}
\end{par}

\begin{par}
\begin{flushleft}
Графики входа и выходов на рис. ???.
\end{flushleft}
\end{par}

\begin{par}
\begin{flushleft}
???РИСУНОК???
\end{flushleft}
\end{par}

\begin{par}
\begin{flushleft}
Графики ошибки на рис. ???.
\end{flushleft}
\end{par}

\begin{par}
\begin{flushleft}
???РИСУНОК???
\end{flushleft}
\end{par}

\begin{par}
\begin{flushleft}
Как мы видим из полученных графиков, при увеличении $k$ уменьшается установившаяся ошибка, наблюдается незначительное повышение колебательности.
\end{flushleft}
\end{par}

\begin{par}
\begin{flushleft}
\textbf{Движение с постоянным ускорением: }$G(s)=\frac{a}{s^3 }$.
\end{flushleft}
\end{par}

\begin{par}
\begin{flushleft}
Выполним моделирование. Структурная схема представлена на рис. ???.
\end{flushleft}
\end{par}

\begin{par}
\begin{flushleft}
???РИСУНОК???
\end{flushleft}
\end{par}

\begin{par}
\begin{flushleft}
Графики входа и выходов на рис. ???.
\end{flushleft}
\end{par}

\begin{par}
\begin{flushleft}
???РИСУНОК???
\end{flushleft}
\end{par}

\begin{par}
\begin{flushleft}
Графики ошибки на рис. ???.
\end{flushleft}
\end{par}

\begin{par}
\begin{flushleft}
???РИСУНОК???
\end{flushleft}
\end{par}

\begin{par}
\begin{flushleft}
Из графиков видно, что следить за таким входным воздействием у интегрального регулятора уже не получается - ошибка растет линейно (с точностью аппроксимации функции ошибки). При увеличении $k$ ошибка накапливается медленнее, но при $t\to \infty$ она всё равно уйдет в бесконечность из-за недостающего порядка астатизма.
\end{flushleft}
\end{par}

\begin{par}
\begin{flushleft}
\textbf{Выводы}
\end{flushleft}
\end{par}

\begin{par}
\begin{flushleft}
В данном задании мы исследовали интегральный регулятор. Для анализа системы, как и ранее, было удобно использовать теорему о конечном значении, а для проверки устойчивости был применен критерий Гурвица, который показал существенное ограничение на максимальный коэффициент усиления $k$.
\end{flushleft}
\end{par}

\begin{par}
\begin{flushleft}
При исследовании стационарного режима было рассчитано, что установившаяся ошибка равна нулю для любого значения коэффициента $k$ из диапазона устойчивости. Этот аналитический результат подтверждается результатами моделирования и объясняется тем, что система с И-регулятором обладает астатизмом первого порядка, что позволяет ей без ошибки отрабатывать постоянное задающее воздействие.
\end{flushleft}
\end{par}

\begin{par}
\begin{flushleft}
При исследовании режима движения с постоянной скоростью были аналитически получены значения установившейся ошибки, которые подтверждаются результатами моделирования. Увеличение коэффициента усиления $k$ приводит к уменьшению установившейся ошибки слежения. Однако вместе с этим может незначительно возрастать колебательность переходного процесса, а максимально допустимое значение $k$ жестко ограничено условиями устойчивости.
\end{flushleft}
\end{par}

\begin{par}
\begin{flushleft}
При исследовании режима движения с постоянным ускорением показано, что интегральный регулятор не обеспечивает нулевой или фиксированной установившейся ошибки. Ошибка увеличивается со временем, а увеличение коэффициента усиления лишь уменьшает скорость её нарастания, но не устраняет её полностью. Это объясняется тем, что система с И-регулятором обладает первым порядком астатизма, следовательно, она способна компенсировать без ошибки постоянные воздействия и с конечной ошибкой — линейно растущие, но не способна отработать квадратично (и быстрее) растущие воздействия.
\end{flushleft}
\end{par}

\begin{par}
\begin{flushleft}
Таким образом, использование интегрального регулятора является эффективным для задач слежения за постоянным сигналом и допустимым для сигналов, растущих с постоянной скоростью, если допустима ненулевая установившаяся ошибка, величина которой может быть уменьшена ростом $k$ в пределах устойчивости. Однако для слежения за ускоренно растущими сигналами данного регулятора недостаточно.
\end{flushleft}
\end{par}

\begin{matlabcode}
k_arr = [ 0.125, 0.07, 0.02 ];
A = 3;
V = 2.5;
half_a = 0.4;

simtime = 400;
figure;
set(gcf,'AutoResizeChildren','off');
hold on;
colors = hsv(length(k_arr));

for i = 1:length(k_arr)
    k = k_arr(i);
    plot_4 = sim("lab4_41.slx", simtime);
    plot(plot_4.y.Time, plot_4.y.Data, LineWidth=2, LineStyle="-", ...
                DisplayName=sprintf('$y(t), k=%.3f$', k), Color=colors(i,:))
    hold on;
end
plot_41 = sim("lab4_41.slx", simtime);
plot(plot_41.g.Time, plot_41.g.Data, LineWidth=3, color="black", LineStyle="--", ...
            DisplayName='$g(t)=3$')
hold on;

grid on;
legend(Location="best", FontSize=14, Interpreter="latex");
xlabel("$t, c$", Interpreter="latex");
ylabel("$f(t)$", Interpreter="latex");
set(gcf, Position=[100, 100, 1200, 600]);   
set(gca, FontSize=14);

drawnow;
exportgraphics(gcf, "images/plot_41.png", Resolution=200)
hold off;
figure;
set(gcf,'AutoResizeChildren','off');
hold on;

for i = 1:length(k_arr)
    k = k_arr(i);
    plot_4e = sim("lab4_41.slx", simtime);
    plot(plot_4e.e.Time, plot_4e.e.Data, LineWidth=2, LineStyle="-", ...
                DisplayName=sprintf('$e(t), k=%.3f$', k))
    hold on;
end

grid on;
legend(Location="best", FontSize=14, Interpreter="latex");
xlabel("$t, c$", Interpreter="latex");
ylabel("$e(t)$", Interpreter="latex");
set(gcf, Position=[150, 150, 1200, 600]);
set(gca, FontSize=14);

drawnow;
exportgraphics(gcf, "images/plot_41_err.png", Resolution=200)
hold off;
\end{matlabcode}


\begin{matlabcode}
simtime = 250;
figure;
set(gcf,'AutoResizeChildren','off');
hold on;

for i = 1:length(k_arr)
    k = k_arr(i);
    plot_4 = sim("lab4_42.slx", simtime);
    plot(plot_4.y.Time, plot_4.y.Data, LineWidth=2, LineStyle="-", ...
                DisplayName=sprintf('$y(t), k=%.3f$', k))
    hold on;
end
plot_42 = sim("lab4_42.slx", simtime);
plot(plot_42.g.Time, plot_42.g.Data, LineWidth=2, color="red", LineStyle="--", ...
            DisplayName='$g(t)=2.5t$')
hold on;

grid on;
legend(Location="best", FontSize=14, Interpreter="latex");
xlabel("$t, c$", Interpreter="latex");
ylabel("$f(t)$", Interpreter="latex");
set(gcf, Position=[100, 100, 1200, 600]);   
set(gca, FontSize=14);

drawnow;
exportgraphics(gcf, "images/plot_42.png", Resolution=200)
hold off;
e_inf4 = [ 20, 35.7, 125]; 

figure;
set(gcf,'AutoResizeChildren','off');
hold on;
colors = hsv(length(k_arr));

for i = 1:length(k_arr)
    k = k_arr(i);
    plot_4e = sim("lab4_42.slx", simtime);
    plot(plot_4e.e.Time, plot_4e.e.Data, LineWidth=2, LineStyle="-", ...
                DisplayName=sprintf('$e(t), k=%.3f$', k), Color=colors(i,:))
    darkColor = colors(i,:) * (1 - 0.5);

    yline(e_inf4(i), '--', Color=darkColor, LineWidth=2, ...
        DisplayName=sprintf('$e_{\\infty}, k=%.3f$', k));

    hold on;
end

grid on;
legend(Location="best", FontSize=14, Interpreter="latex");
xlabel("$t, c$", Interpreter="latex");
ylabel("$e(t)$", Interpreter="latex");
set(gcf, Position=[150, 150, 1200, 600]);
set(gca, FontSize=14);

drawnow;
exportgraphics(gcf, "images/plot_42_err.png", Resolution=200)
hold off;
\end{matlabcode}


\begin{matlabcode}
simtime = 400;
figure;
set(gcf,'AutoResizeChildren','off');
hold on;

for i = 1:length(k_arr)
    k = k_arr(i);
    plot_4 = sim("lab4_43.slx", simtime);
    plot(plot_4.y.Time, plot_4.y.Data, LineWidth=2, LineStyle="-", ...
                DisplayName=sprintf('$y(t), k=%.3f$', k))
    hold on;
end
plot_43 = sim("lab4_43.slx", simtime);
plot(plot_43.g.Time, plot_43.g.Data, LineWidth=2, color="red", LineStyle="--", ...
            DisplayName='$g(t)=0.4t^2$')
hold on;

grid on;
legend(Location="best", FontSize=14, Interpreter="latex");
xlabel("$t, c$", Interpreter="latex");
ylabel("$f(t)$", Interpreter="latex");
set(gcf, Position=[100, 100, 1200, 600]);   
set(gca, FontSize=14);

drawnow;
exportgraphics(gcf, "images/plot_43.png", Resolution=200)
hold off;
figure;
set(gcf,'AutoResizeChildren','off');
hold on;

for i = 1:length(k_arr)
    k = k_arr(i);
    plot_4e = sim("lab4_43.slx", simtime);
    plot(plot_4e.e.Time, plot_4e.e.Data, LineWidth=2, LineStyle="-", ...
                DisplayName=sprintf('$e(t), k=%.3f$', k))
    hold on;
end

grid on;
legend(Location="best", FontSize=14, Interpreter="latex");
xlabel("$t, c$", Interpreter="latex");
ylabel("$e(t)$", Interpreter="latex");
set(gcf, Position=[150, 150, 1200, 600]);
set(gca, FontSize=14);

drawnow;
exportgraphics(gcf, "images/plot_43_err.png", Resolution=200)
hold off;
\end{matlabcode}


\matlabtitle{Задание 5}

\begin{par}
\begin{flushleft}
\textbackslash{}section\{Задача слежения для системы с астатизмом первого порядка (ПИ-регулятор)\}
\end{flushleft}
\end{par}

\begin{par}
\begin{flushleft}
\texttt{Рассмотрим замкнутую систему, заданную структурной схемой на рисунке ???ВЫШЕ???, где}
\end{flushleft}
\end{par}

\begin{par}
$$H(s)=\frac{k_{\textrm{И}} }{s}+k_{\textrm{П}}$$
\end{par}

\begin{par}
\begin{flushleft}
\texttt{- пропорционально-интегральный регулятор. Варианты передаточной функции объекта управления }$W(s)$\texttt{ и характеристики задающего воздействия }$g(t)$\texttt{ приведены в }\textbf{Таблице 1 }инструкции по выполнению лабораторной работы\texttt{. Зададимся не менее чем двумя значениями каждого из параметров }$k_{\textrm{И}}$\texttt{ и }$k_{\textrm{П}}$\texttt{ и составим все возможные комбинации пар выбранных значений параметров (т.е. не менее четырёх комбинаций). Все пары параметров должны обеспечивать асимптотически устойчивую замкнутую систему.}
\end{flushleft}
\end{par}

\begin{par}
\begin{flushleft}
\texttt{Исследуем режим движения с постоянной скоростью при }$g(t)=Vt$\texttt{: выполним моделирование для выбранных значений }$k$\texttt{, а также аналитически определим предельное значение ошибки }$e_{уст}$\texttt{ для каждого значения }$k$\texttt{. Сопоставим результаты между собой и сделаем выводы.}
\end{flushleft}
\end{par}

\begin{par}
\begin{flushleft}
\texttt{Исследуем режим слежения за гармоническим сигналом при }$g(t)=a\sin (\omega t)$\texttt{: выполним моделирование для выбранных значений }$k$\texttt{. Сопоставим результаты между собой и сделаем выводы.}
\end{flushleft}
\end{par}


\vspace{1em}
\begin{par}
\begin{flushleft}
Возьмём из таблицы для варианта 27 передаточную функцию и параметры управляющего воздействия:
\end{flushleft}
\end{par}

\begin{par}
$$W(s)=\frac{1}{4s^2 +0.6s+1},~~V=2.5,~~A\sin (\omega t)=3\sin (0.4t)$$
\end{par}

\begin{par}
\begin{flushleft}
Передаточная функция замкнутой системы:
\end{flushleft}
\end{par}

\begin{par}
$$\begin{array}{l}
W_{g\to y} (s)=\frac{W(s)H(s)}{1+W(s)H(s)}=\frac{\left(k_П +\frac{k_И }{s}\right)\cdot \frac{1}{4s^2 +0.6s+1}}{1+\left(k_П +\frac{k_И }{s}\right)\cdot \frac{1}{4s^2 +0.6s+1}}=\frac{k_П s+k_И }{s(4s^2 +0.6s+1)+k_П s+k_И }=\\
=\frac{k_П s+k_И }{4s^3 +0.6s^2 +(1+k_П )s+k_И }
\end{array}$$
\end{par}

\begin{par}
\begin{flushleft}
Найдем такие $k$, при которых замкнутая система является ассимптотически устойчивой.
\end{flushleft}
\end{par}

\begin{par}
\begin{flushleft}
$4s^3 +0.6s^2 +(1+k_П )s+k_И =0$, воспользуемся критерием Гурвица, матрица Гурвица:
\end{flushleft}
\end{par}

\begin{par}
$$A=\left(\begin{array}{ccc}
0.6 & k_И  & 0\\
4 & 1+k_П  & 0\\
0 & 0.6 & k_И 
\end{array}\right)$$
\end{par}

\begin{par}
\begin{flushleft}
Ведущие угловые миноры матрицы Гурвица: $\Delta_1 =0.6,~~\Delta_2 =0.6(1+k_П )-4k_И ,~~\Delta_3 =k_И \cdot \Delta_2 =k_И (0.6(1+k_П )-4k_И )$
\end{flushleft}
\end{par}

\begin{par}
\begin{flushleft}
Следовательно, система ассимптотически устойчива при:
\end{flushleft}
\end{par}

\begin{par}
$$\left\lbrace \begin{array}{c}
k_И >0\\
~~0.6(1+k_П )-4k_И >0\\
k_П >-1
\end{array}\right.~~\Rightarrow \left\lbrace \begin{array}{c}
k_И \in \left(0,~0.15(1+k_П )\right)\\
k_П >-1
\end{array}\right.$$
\end{par}

\begin{par}
\begin{flushleft}
Зададимся значениями $k_П =\lbrace 1,2\rbrace$\texttt{, }$k_И =\lbrace 0.1,0.2\rbrace$\texttt{. }Таким образом, рассмотрим комбинации параметров: $(k_П ,k_И )=(1,0.1),~~(1,0.2),~~(2,0.1),~~(2,0.2)$.
\end{flushleft}
\end{par}

\begin{par}
\begin{flushleft}
Проведем исследование для двух режимов работы.
\end{flushleft}
\end{par}

\begin{par}
\begin{flushleft}
\textbf{Движение с постоянной скоростью: }$G(s)=\frac{V}{s^2 }$.
\end{flushleft}
\end{par}

\begin{par}
\begin{flushleft}
Аналитически определим предельное значение ошибки $e_{уст}$ для каждого значения $k$:
\end{flushleft}
\end{par}

\begin{par}
\begin{flushleft}
Ошибка слежения:
\end{flushleft}
\end{par}

\begin{par}
$$\begin{array}{l}
E(s)=\frac{1}{1+W(s)H(s)}G(s)=\frac{1}{1+\left(k_П +\frac{k_И }{s}\right)\cdot \frac{1}{4s^2 +0.6s+1}}\cdot \frac{V}{s^2 }=\frac{s(4s^2 +0.6s+1)}{s(4s^2 +0.6s+1)+k_П s+k_И }\cdot \frac{V}{s^2 }=\\
=\frac{V(4s^2 +0.6s+1)}{s(4s^3 +0.6s^2 +(1+k_П )s+k_И )}
\end{array}$$
\end{par}

\begin{par}
\begin{flushleft}
Согласно теореме о конечном значении, если все полюса $s\cdot E(s)$ имеют отрицательную вещественную часть, $\lim_{t\to +\infty } y(t)=\lim_{s\to 0} sY(s)$.
\end{flushleft}
\end{par}

\begin{par}
\begin{flushleft}
$s\cdot E(s)=\frac{V(4s^2 +0.6s+1)}{4s^3 +0.6s^2 +(1+k_П )s+k_И }$, полюса данной передаточной функции отрицательны, так как ранее мы уже подбирали диапазон $k$, соответствующий отрицательным полюсам передаточной функции с таким же знаменателем. Можем применнить теорему о конечном значении и найти установившуюся ошибку:
\end{flushleft}
\end{par}

\begin{par}
$$e_{\textrm{уст}} =\lim_{s\to 0} s\cdot E(s)=\lim_{s\to 0} \frac{V(4s^2 +0.6s+1)}{4s^3 +0.6s^2 +(1+k_П )s+k_И }=\frac{V(4\cdot 0^2 +0.6\cdot 0+1)}{4\cdot 0^3 +0.6\cdot 0^2 +(1+k_П )\cdot 0+k_И }=\frac{V}{k_И }$$
\end{par}

\begin{par}
\begin{flushleft}
Как видим, пропорциональный коэффициент не влияет на значение установившейся ошибки.
\end{flushleft}
\end{par}

\begin{par}
\begin{flushleft}
Рассчитаем для каждого $k_И$:
\end{flushleft}
\end{par}

\begin{itemize}
\setlength{\itemsep}{-1ex}
   \item{\begin{flushleft} $\displaystyle k_И =0.1,~~e_{\textrm{уст}} =\frac{2.5}{0.1}=25$ \end{flushleft}}
   \item{\begin{flushleft} $\displaystyle k_И =0.2,~~e_{\textrm{уст}} =\frac{2.5}{0.2}=12.5$ \end{flushleft}}
\end{itemize}

\begin{par}
\begin{flushleft}
Выполним моделирование. Структурная схема представлена на рис. ???.
\end{flushleft}
\end{par}

\begin{par}
\begin{flushleft}
???РИСУНОК???
\end{flushleft}
\end{par}

\begin{par}
\begin{flushleft}
Графики входа и выходов на рис. ???.
\end{flushleft}
\end{par}

\begin{par}
\begin{flushleft}
???РИСУНОК???
\end{flushleft}
\end{par}

\begin{par}
\begin{flushleft}
Графики ошибки на рис. ???.
\end{flushleft}
\end{par}

\begin{par}
\begin{flushleft}
???РИСУНОК???
\end{flushleft}
\end{par}

\begin{par}
\begin{flushleft}
По графикам видим экспериментальное потверждение того, что пропорциональный коэффициент не влияет на установившуюся ошибку, при повышении интегрального коэффициента у нас снова получается уменьшение ошибки и небольшое увеличение колебаний, но идеально уследить за сигналом всё ещё не даёт порядок астатизма.
\end{flushleft}
\end{par}

\begin{par}
\begin{flushleft}
\textbf{Слежение за гармоническим сигналом: }$G(s)=\frac{a\omega }{s^2 +\omega^2 }$.
\end{flushleft}
\end{par}

\begin{par}
\begin{flushleft}
Выполним моделирование. Структурная схема представлена на рис. ???.
\end{flushleft}
\end{par}

\begin{par}
\begin{flushleft}
???РИСУНОК???
\end{flushleft}
\end{par}

\begin{par}
\begin{flushleft}
Графики входа и выходов на рис. ???.
\end{flushleft}
\end{par}

\begin{par}
\begin{flushleft}
???РИСУНОК???
\end{flushleft}
\end{par}

\begin{par}
\begin{flushleft}
Графики ошибки на рис. ???.
\end{flushleft}
\end{par}

\begin{par}
\begin{flushleft}
???РИСУНОК???
\end{flushleft}
\end{par}

\begin{par}
\begin{flushleft}
В данном случае видим, что ошибка не сходится к определённому значению, а становится синусоидальной, но с установишейся амплитудой и частотой. Идеально проследить за сигналом снова не получается, но увеличение интегрального коэффициента делает амплитуду ошибку меньше. П-коэффициент влияет очень слабо.
\end{flushleft}
\end{par}

\begin{par}
\begin{flushleft}
\textbf{Выводы}
\end{flushleft}
\end{par}

\begin{par}
\begin{flushleft}
В данном задании мы исследовали пропорционально-интегральный регулятор. Были найдены подходящие для ассимптотической устойчивости диапазоны коэффициентов, их зависимость между собой.
\end{flushleft}
\end{par}

\begin{par}
\begin{flushleft}
При исследовании режима движения с постоянной скоростью мы аналитически нашли значения установившейся ошибки и показали, что П-коэффициент не влияет на неё. При моделировании данные результаты были подтверждены экспериментом. Увеличение интегрального коэффициента приводит к уменьшению установившейся ошибки, пропорциональный коэффициент влияет слабо. Так же отметим, что данные результаты хорошо согласутся с выводами из предыдущих заданий.
\end{flushleft}
\end{par}

\begin{par}
\begin{flushleft}
При исследовании гармонического сигнала мы обнаружили, что не получается уследить за ним идеально, график ошибки является синусоидальным, и при увеличении И-коэффициента уменьшается его амплитуда, П-коэффициент так же влияет очень слабо.
\end{flushleft}
\end{par}

\begin{par}
\begin{flushleft}
Мы получили, и тем самым подтвердили теоретические ожидания, что ПИ-регулятор сочетает в себе свойства П и И-регуляторов. Пропорциональный коэффициент играет маленькую роль в нём при работе в тех режимах, которые мы исследовали в данном задании. Мы можем увеличивать И-коэффициент, чтобы уменьшать ошибку при слежении за линейно растущим и синусоидальным входным воздействием, но не сможем их полностью "догнать".
\end{flushleft}
\end{par}

\begin{matlabcode}
k_arr = [1, 0.1; 2, 0.1; 1, 0.2; 2, 0.2];
A = 3;
V = 2.5;
w = 0.4;

simtime = 80;
figure;
set(gcf,'AutoResizeChildren','off');
hold on;
ls = ["-", "--", "-", "--"];

for i = 1:size(k_arr,1)
    k_p = k_arr(i,1);
    k_i = k_arr(i,2);

    plot_5 = sim("lab4_51.slx", simtime);
    plot(plot_5.y.Time, plot_5.y.Data, LineWidth=3, LineStyle=ls(i), ...
                DisplayName=sprintf('$y(t), k_p=%.1f, k_i=%.1f$', k_p, k_i))
    hold on;
end
plot_51 = sim("lab4_51.slx", simtime);
plot(plot_51.g.Time, plot_51.g.Data, LineWidth=2, color="red", LineStyle="--", ...
            DisplayName='$g(t)=2.5t$')
hold on;

grid on;
legend(Location="best", FontSize=14, Interpreter="latex");
xlabel("$t, c$", Interpreter="latex");
ylabel("$f(t)$", Interpreter="latex");
set(gcf, Position=[100, 100, 1200, 600]);   
set(gca, FontSize=14);

drawnow;
exportgraphics(gcf, "images/plot_51.png", Resolution=200)
hold off;
\end{matlabcode}


\begin{matlabcode}
e_inf5 = [ 25, 12.5]; 
ki = [ 0.1, 0.2 ];
simtime = 180;
figure;
set(gcf,'AutoResizeChildren','off');
hold on;
colors = hsv(length(k_arr));

for i = 1:2
    darkColor = colors(i,:) * (1 - 0.5);
    yline(e_inf5(i), '--', Color=darkColor, LineWidth=2, ...
        DisplayName=sprintf('$e_{\\infty}, k_i=%.1f$', ki(i)));

    hold on;
end

for i = 1:size(k_arr,1)
    k_p = k_arr(i,1);
    k_i = k_arr(i,2);
    plot_5e = sim("lab4_51.slx", simtime);
    plot(plot_5e.e.Time, plot_5e.e.Data, LineWidth=2, LineStyle="-", ...
                DisplayName=sprintf('$e(t), k_p=%.1f, k_i=%.1f$', k_p, k_i))

    hold on;
end

grid on;
legend(Location="best", FontSize=14, Interpreter="latex");
xlabel("$t, c$", Interpreter="latex");
ylabel("$e(t)$", Interpreter="latex");
set(gcf, Position=[150, 150, 1200, 600]);
set(gca, FontSize=14);

drawnow;
exportgraphics(gcf, "images/plot_51_err.png", Resolution=200)
hold off;
\end{matlabcode}


\begin{matlabcode}
k_arr = [1, 0.1; 2, 0.1; 1, 0.2; 2, 0.2];
A = 3;
V = 2.5;
w = 0.4;
simtime = 60;
figure;
set(gcf,'AutoResizeChildren','off');
hold on;

for i = 1:size(k_arr,1)
    k_p = k_arr(i,1);
    k_i = k_arr(i,2);
    plot_5 = sim("lab4_52.slx", simtime);
    plot(plot_5.y.Time, plot_5.y.Data, LineWidth=2, LineStyle="-", ...
                DisplayName=sprintf('$y(t), k_p=%.1f, k_i=%.1f$', k_p, k_i))
    hold on;
end
plot_52 = sim("lab4_52.slx", simtime);
plot(plot_52.g.Time, plot_52.g.Data, LineWidth=2, color="black", LineStyle="--", ...
            DisplayName='$g(t)=3\sin(0.4t)$')
hold on;

grid on;
legend(Location="southoutside", FontSize=14, Interpreter="latex");
xlabel("$t, c$", Interpreter="latex");
ylabel("$f(t)$", Interpreter="latex");
set(gcf, Position=[100, 100, 1200, 800]);   
set(gca, FontSize=14);

drawnow;
exportgraphics(gcf, "images/plot_52.png", Resolution=200)
hold off;
\end{matlabcode}


\begin{matlabcode}
simtime = 150;
figure;
set(gcf,'AutoResizeChildren','off');
hold on;
colors = hsv(length(k_arr));

ls = ["-", "-", "-", "--"];
for i = 1:size(k_arr,1)
    k_p = k_arr(i,1);
    k_i = k_arr(i,2);
    plot_5e = sim("lab4_52.slx", simtime);
    plot(plot_5e.e.Time, plot_5e.e.Data, LineWidth=3, LineStyle=ls(i), ...
                DisplayName=sprintf('$e(t), k_p=%.1f, k_i=%.1f$', k_p, k_i))

    hold on;
end

grid on;
legend(Location="southoutside", FontSize=14, Interpreter="latex");
xlabel("$t, c$", Interpreter="latex");
ylabel("$e(t)$", Interpreter="latex");
set(gcf, Position=[150, 150, 1200, 800]);
set(gca, FontSize=14);

drawnow;
exportgraphics(gcf, "images/plot_52_err.png", Resolution=200)
hold off;
\end{matlabcode}


\matlabtitle{Задание 6}

\begin{par}
\begin{flushleft}
\textbackslash{}section\{Задача слежения за гармоническим сигналом (регулятор общего вида)\}
\end{flushleft}
\end{par}

\begin{par}
\begin{flushleft}
\texttt{Рассмотрим замкнутую систему, заданную структурной схемой на рисунке ???ВЫШЕ???, где}
\end{flushleft}
\end{par}

\begin{par}
$$H(s)=\frac{\sum_{k=0}^m (b_k s^k )}{\sum_{k=0}^m (a_k s^k )}~~(4)$$
\end{par}

\begin{par}
\begin{flushleft}
\texttt{— регулятор общего вида. Варианты передаточной функции объекта управления }$W(s)$\texttt{ и характеристики задающего воздействия }$g(t)$\texttt{ приведены в }\textbf{Таблице 1 }инструкции по выполнению лабораторной работы\texttt{.}
\end{flushleft}
\end{par}

\begin{par}
\begin{flushleft}
\texttt{Для задающего гармонического сигнала }$g(t)=A\sin (\omega t)$\texttt{ синтезируем физически реализуемый регулятор вида (4), способный обеспечить предельное значение ошибки }$e_{\textrm{уст}} =0$\texttt{:}
\end{flushleft}
\end{par}

\begin{itemize}
\setlength{\itemsep}{-1ex}
   \item{\begin{flushleft} \texttt{Определим необходимый порядок }$m$; \end{flushleft}}
   \item{\begin{flushleft} \texttt{Зададимся типовым полиномом Ньютона соответствующего порядка;} \end{flushleft}}
   \item{\begin{flushleft} \texttt{Зададимся конкретным значением коэффициента подобия }$\omega_0$; \end{flushleft}}
   \item{\begin{flushleft} \texttt{Вычислим конкретные параметры }$b_k$\texttt{ и }$a_k$; \end{flushleft}}
   \item{\begin{flushleft} \texttt{Подтвердим работоспособность аналитическим расчетом }$e_{\textrm{уст}}$; \end{flushleft}}
   \item{\begin{flushleft} \texttt{Выполним моделирование, демонстрирующее работоспособность синтезированного регулятора;} \end{flushleft}}
   \item{\begin{flushleft} \texttt{Сравним фактическое время переходного процесса для ошибки }$e(t)$ \texttt{с теоретическим, полученным на основании }\texttt{\textbf{Теоремы подобия}}\texttt{.    } \end{flushleft}}
\end{itemize}

\begin{par}
\begin{flushleft}
Возьмём из таблицы для варианта 27 передаточную функцию и параметры управляющего воздействия:
\end{flushleft}
\end{par}

\begin{par}
$$W(s)=\frac{1}{4s^2 +0.6s+1},~~A\sin (\omega t)=3\sin (0.4t)$$
\end{par}

\begin{par}
\begin{flushleft}
Запишем на регулятор в виде $H(s)=\frac{R(s)}{Q(s)}$. Передаточная функция по ошибке:
\end{flushleft}
\end{par}

\begin{par}
$$\begin{array}{l}
W_{G\to E} (s)=\frac{1}{1+W(s)H(s)}=\frac{1}{1+W(s)\frac{R(s)}{Q(s)}}=\frac{Q(s)}{Q(s)+W(s)R(s)}=\\
\frac{Q(s)(4s^2 +0.6s+1)}{Q(s)(4s^2 +0.6s+1)+R(s)}
\end{array}$$
\end{par}

\begin{par}
\begin{flushleft}
Образ Лапласа для нашего сигнала: $G(s)=\frac{a\omega }{s^2 +\omega^2 }$
\end{flushleft}
\end{par}

\begin{par}
\begin{flushleft}
Попробуем применить теорему о конечном значении:
\end{flushleft}
\end{par}

\begin{par}
$$\begin{array}{l}
\lim_{t\to \infty } e(t)=\lim_{s\to 0} \left(sW_{G\to E} (s)G(s)\right)=\lim_{s\to 0} \left(\frac{Q(s)(4s^2 +0.6s+1)}{Q(s)(4s^2 +0.6s+1)+R(s)}\cdot \frac{3s\omega }{s^2 +0.4^2 }\right)=\\
\lim_{s\to 0} \left(\frac{Q(s)(4s^2 +0.6s+1)3s\omega }{(Q(s)(4s^2 +0.6s+1)+R(s))(s^2 +0.4^2 )}\right)
\end{array}$$
\end{par}

\begin{par}
\begin{flushleft}
\texttt{Чтобы мы могли её применить, возьмём }$Q(s)=Q^{\prime } (s)(s^2 +0.4^2 )$:
\end{flushleft}
\end{par}

\begin{par}
$$\lim_{t\to \infty } e(t)=\lim_{s\to 0} \left(\frac{Q^{\prime } (s)(4s^2 +0.6s+1)3s\omega }{Q^{\prime } (s)(s^2 +0.16)(4s^2 +0.6s+1)+R(s)}\right)=0$$
\end{par}

\begin{par}
\begin{flushleft}
Но у $Q^{\prime } (s)(s^2 +0.16)(4s^2 +0.6s+1)+R(s)$ должны быть только корни с отрицательной вещ. частью.
\end{flushleft}
\end{par}

\begin{par}
\begin{flushleft}
\texttt{Пусть }$Q^{\prime } (s)=a_2 =const$\texttt{, тогда }$Q(s)=a_2 s^2 +0.16a_2$. \texttt{Пусть }$R(s)=b_2 s^2 +b_1 s+b_0$\texttt{, тогда:}
\end{flushleft}
\end{par}

\begin{par}
\begin{flushleft}
$(a_2 s^2 +0.16a_2 )(4s^2 +0.6s+1)+(b_2 s^2 +b_1 s+b_0 )=4a_2 s^4 +0.6a_2 s^3 +(1.64a_2 +b_2 )s^2 +(0.096a_2 +b_1 )s+(0.16a_2 +b_0 )$;
\end{flushleft}
\end{par}

\begin{par}
$$\left\lbrace \begin{array}{c}
c_0 =0.16a_2 +b_0 \\
c_1 =0.096a_2 +b_1 \\
c_2 =1.64a_2 +b_2 \\
c_3 =0.6a_2 \\
c_4 =4a_2 
\end{array}\right.$$
\end{par}

\begin{par}
\begin{flushleft}
Порядок регулятора $m=2$.
\end{flushleft}
\end{par}

\begin{par}
\begin{flushleft}
Порядок замкнутой системы 4, типовой полином Ньютона: $s^4 +4\omega_0 s^3 +6\omega_0^2 s^2 +4\omega_0^3 s+\omega_0^4 .$
\end{flushleft}
\end{par}

\begin{par}
\begin{flushleft}
Найдем коэффициент подобия. Поделим все коэффициенты на $4a_2$:
\end{flushleft}
\end{par}

\begin{par}
$$\left\lbrace \begin{array}{c}
4\omega_0 =0.15~~\Rightarrow ~~\omega_0 =\frac{3}{80},\\
6\omega_0^2 =0.41+\frac{b_2 }{4a_2 },\\
4\omega_0^3 =0.024+\frac{b_1 }{4a_2 },\\
\omega_0^4 =0.04+\frac{b_0 }{4a_2 }.
\end{array}\right.$$
\end{par}

\begin{par}
$$\omega_0^2 =\frac{9}{6400},~~\omega_0^3 =\frac{27}{512000},~~\omega_0^4 =\frac{81}{40960000}$$
\end{par}

\begin{par}
\begin{flushleft}
Вычислим параметры:
\end{flushleft}
\end{par}

\begin{par}
$$\left\lbrace \begin{array}{c}
\frac{b_2 }{4a_2 }=6\cdot \frac{9}{6400}-0.41=\frac{54}{6400}-\frac{2624}{6400}=-\frac{2570}{6400}=-\frac{1285}{3200},\\
\frac{b_1 }{4a_2 }=4\cdot \frac{27}{512000}-0.024=\frac{108}{512000}-\frac{12288}{512000}=-\frac{12180}{512000}=-\frac{609}{25600},\\
\frac{b_0 }{4a_2 }=\frac{81}{40960000}-0.04=\frac{81}{40960000}-\frac{1638400}{40960000}=-\frac{1638319}{40960000}
\end{array}\right.$$
\end{par}

\begin{par}
\begin{flushleft}
Возьмём $a_2 =10240000$:
\end{flushleft}
\end{par}

\begin{par}
$$b_2 =-1285\cdot 12800=-16448000,~~b_1 =-609\cdot 1600=-974400,~~b_0 =-1638319.$$
\end{par}

\begin{par}
\begin{flushleft}
Передаточная функция регулятора: $H(s)=\frac{-16448000s^2 -974400s-1638319}{10240000s^2 +1638400}$
\end{flushleft}
\end{par}

\begin{par}
\begin{flushleft}
\texttt{Подтвердим работоспособность аналитическим расчетом }$e_{\textrm{уст}}$\texttt{:}
\end{flushleft}
\end{par}

\begin{par}
\begin{flushleft}
\texttt{С учетом того, что полином Ньютона имеет вид }$(s+\omega_0 )^4$ (для степени 4), $\omega_0 >0$, полюса $s\cdot E(s)$ отрицательны и мы можем применить теорему о конечном значении:
\end{flushleft}
\end{par}

\begin{par}
$$\lim_{t\to \infty } e(t)=\lim_{s\to 0} \left(\frac{3\omega a_2 (4s^3 +0.6s^2 +s)}{c_4 s^4 +c_3 s^3 +c_2 s^2 +c_1 s+c_0 }\right)=\frac{3\cdot 0.4\cdot 10240000\cdot (4\cdot 0+0.6\cdot 0+0)}{0+0+0+0+c_0 }=\frac{0}{c_0 }=0$$
\end{par}

\begin{par}
\begin{flushleft}
Система работоспособна.
\end{flushleft}
\end{par}

\begin{par}
\begin{flushleft}
Выполним моделирование. Структурная схема представлена на рис. ???.
\end{flushleft}
\end{par}

\begin{par}
\begin{flushleft}
???РИСУНОК???
\end{flushleft}
\end{par}

\begin{par}
\begin{flushleft}
Графики входа и выходов на рис. ???.
\end{flushleft}
\end{par}

\begin{par}
\begin{flushleft}
???РИСУНОК???
\end{flushleft}
\end{par}

\begin{par}
\begin{flushleft}
Графики ошибки на рис. ???.
\end{flushleft}
\end{par}

\begin{par}
\begin{flushleft}
???РИСУНОК???
\end{flushleft}
\end{par}

\begin{par}
\begin{flushleft}
Время переходного процесса на основании \textbf{Теоремы подобия}: $t_{\textrm{П}}^{\prime } =\frac{t_{\textrm{П}} }{\omega_0 }=\frac{7.8}{3/80}=208$[c].
\end{flushleft}
\end{par}

\begin{par}
\begin{flushleft}
Результат, полученный экспериментально для $e_{\textrm{уст}} <0.05$: $t_{\textrm{П,}\;\textrm{эксп}}^{\prime } =465.958$[c].
\end{flushleft}
\end{par}

\begin{par}
\begin{flushleft}
\textbf{Выводы}
\end{flushleft}
\end{par}


\vspace{1em}

\vspace{1em}
\begin{par}
\begin{flushleft}
\textbf{Общие выводы}
\end{flushleft}
\end{par}


\vspace{1em}
\begin{matlabcode}
simtime = 500;
figure;
set(gcf,'AutoResizeChildren','off');
hold on;
plot_6 = sim("lab4_6.slx", simtime);
plot(plot_6.y.Time, plot_6.y.Data, LineWidth=3, LineStyle="-", ...
            DisplayName="$y(t)$")
hold on;
plot_61 = sim("lab4_6.slx", simtime);
plot(plot_61.g.Time, plot_61.g.Data, LineWidth=2, color="red", LineStyle="-", ...
            DisplayName='$g(t)=3\sin(0.4t)$')
hold on;

grid on;
legend(Location="best", FontSize=14, Interpreter="latex");
xlabel("$t, c$", Interpreter="latex");
ylabel("$f(t)$", Interpreter="latex");
set(gcf, Position=[100, 100, 1200, 800]);
set(gca, FontSize=14);

drawnow;
exportgraphics(gcf, "images/plot_6.png", Resolution=200)
hold off;
figure;
set(gcf,'AutoResizeChildren','off');
hold on;

plot(plot_6.e.Time, plot_6.e.Data, LineWidth=3, ...
            DisplayName="$e(t)$")

grid on;
legend(Location="best", FontSize=14, Interpreter="latex");
xlabel("$t, c$", Interpreter="latex");
ylabel("$e(t)$", Interpreter="latex");
set(gcf, Position=[150, 150, 1200, 800]);
set(gca, FontSize=14);

drawnow;
exportgraphics(gcf, "images/plot_6_err.png", Resolution=200)
hold off;
g = plot_6.g.Data;
t = plot_6.g.Time;

e = plot_6.e.Data;

idx = find(abs(e) <= 0.05 * abs(g));

settlingIdx = idx(find(idx > 0.9*length(idx),1));

if ~isempty(settlingIdx)
    settlingTime = t(settlingIdx);
else
    settlingTime = NaN;
end
disp("Время переходного процесса для e(t)<0.05: " + settlingTime + " с");
\end{matlabcode}

\end{document}
